\section{Conclusion}
\label{sec:conclusion}

We characterize the inference cost of block-diffusion (DG) versus
autoregressive (AR) decoding for structured data extraction.
Block-diffusion has a staircase cost shape, unlike the linear cost of
AR. The taxi fare climbs with every token; the subway charges per
station. The shapes cross at $L^{*} = \crossoverL$ output tokens, above
which DG is faster. The advantage is moderate: two and a half times at
natural output, narrowing to just over parity at a thousand tokens. In
cost terms, that is \$1{,}200 per million events, or \$2{,}600 a year on
one GPU at ten thousand events a day. A larger AR does not close it.
Two things do: batching, and co-residency with a model small enough to
fit many copies on one GPU.

Against a MoE-matched AR, this is a fair comparison, and the match is
what makes the crossover interesting. The dense AR-3B decodes at
6.3\,ms/token against the MoE AR-30B's \arSlope\,ms/token. Moving to
the matched baseline adds only 15\% to the serial per-token cost, and
the crossover stays at 48 tokens, close to the dense AR-3B's 59. DG
wins at every practical output length in serial, though the margin is
narrower than an earlier draft reported.

The advantage is limited where batching is available and where tail
SLAs are tight. But in the batch-1 regime---serial investigation
chains, autonomous driving, algorithmic trading---the staircase shape
provides a structural latency advantage that AR cannot match. The $K$
parameter is DG's degradation mechanism: when the deadline tightens, DG
degrades gracefully, producing lower quality with a complete canvas.
AR's only response is truncation. Block-diffusion is the right choice
when output is long, batch-1 is required, and schema completeness
matters. We name five conditions and ground them in concrete workloads.
